% =====================================================================
\clearpage
\part{The reduction}
% =====================================================================
\section{Introduction}

\label{sec:intro}
% =====================================================================

\subsection{The Clay statement, exactly}

The Clay problem, as posed by Jaffe and Witten, is:

\begin{qbox}
Prove that for any compact simple gauge group $G$, a non-trivial quantum Yang--Mills theory
exists on $\mathbb{R}^4$ and has a mass gap $\Delta>0$. Existence includes establishing
axiomatic properties at least as strong as those cited in Streater--Wightman (1964) and
Osterwalder--Schrader (1973, 1975). [Jaffe--Witten, \emph{Quantum Yang--Mills Theory}, Clay
Mathematics Institute]
\end{qbox}

Three features of this statement are habitually under-weighted and drive everything that follows.

\textbf{(1) Uniformity.} It demands a uniform gap: not merely that each finite lattice volume has
a positive spectral gap (trivially true by compactness), but that the gap survives $L\to\infty$
and $a\to0$ with a volume- and cutoff-uniform lower bound, and that the resulting theory clusters
exponentially at a rate $m>0$ in every channel carrying the lightest state.

\textbf{(2) Gauge invariance.} It demands that the gap be a feature of the gauge-invariant,
physical theory --- not of a gauge-fixed chart, and not of a perturbative expansion around a
chosen background. This sounds like a technicality; it is not. Several of the campaign's
apparent ``positive signals'' turned out to live on gauge-dependent modes, and one of the
no-gos below is precisely that gauge invariance reaches the vacuum but not the first excited
scalar.

\textbf{(3) Existence plus gap are logically separate.} A construction of the measure with the
Osterwalder--Schrader axioms (existence) is a different task from the lower bound on the
Hamiltonian (gap). They interact --- a uniform gap gives clustering, and clustering is one of the
axioms --- but neither implies the other. The campaign treats them as two halves, G1 and G2, and
is explicit that progress on one does not constitute progress on the other.

\figthreehalves

The campaign's operational form of the Clay conclusion, assembled and partially machine-verified
in \code{docs/conditional-theorem.md}, is:
\begin{align*}
&\exists\,\mu\ \text{on}\ \mathbb{R}^4,\ \text{gauge-invariant, satisfying OS0--OS4};\\
&\text{non-Gaussian},\quad g_{\mathrm{ren}}/\Lambda>0;\\
&\Delta/\Lambda\ \ge\ \min(\delta_W,\delta_S)>0;\\
&|S(f\,\tau_a g)-S(f)\,S(g)|\ \le\ C\,e^{-m\|a\|},\qquad
  m\ge c\,\Delta>0,\ \ \text{all gauge-invariant local } f,g;\\
&\text{unique vacuum};\quad \text{area law } \sigma=c_\sigma\Lambda^2>0;\quad
  \text{all constants group-uniform in } (h^{\vee},\dim G,\dim\mathfrak{g}).
\end{align*}

The transfer-matrix form of the same statement is the one used throughout this paper. For the
lattice transfer matrix $T_a=e^{-H_a}$ on the gauge-invariant Hilbert space $\mathcal{H}_a$, with
vacuum $\Omega$,
\begin{equation}
\exists\,\delta>0:\qquad
\inf\operatorname{spec}\!\bigl(H_a\big|_{\Omega^\perp}\bigr)\ \ge\ \delta\Lambda
\qquad\text{for all sufficiently small } a .
\label{eq:clay-transfer}
\end{equation}
The operator $T_a$ is a positive, positivity-improving operator on $\mathcal{H}_a$ at every
$\beta>0$ (the transfer-matrix theorem of \S\ref{thm:pf}); the Clay content is that its first eigenvalue gap does not
collapse as $a\to0$ and $L\to\infty$.

\subsection{The lock-core rewriting}

The campaign works with a sharper, channel-specific rewriting of the gap which, while not
literally equivalent to the Clay statement, is the correct object to attack. Let $\nu$ be the
K\"all\'en--Lehmann spectral measure of the gauge-invariant local operator
$O=\operatorname{tr}F^2$ (more precisely $O_S=\tfrac14\operatorname{tr}F_{\mu\nu}F^{\mu\nu}$,
the scalar density; the normalization is immaterial). Reflection positivity gives $\nu\ge0$
and $\operatorname{supp}\nu\subset[m_1^2,\infty)$, where $m_1$ is the lightest mass in the
$0^{++}$ channel. The \textbf{atom} is
\begin{equation}
f_1^2 \;:=\; \nu(\{m_1^2\}) \;=\; \bigl|\langle 0|\operatorname{tr}F^2|\mathbb{G}\rangle\bigr|^2 ,
\label{eq:atom}
\end{equation}
the squared overlap of the vacuum with the lightest $0^{++}$ glueball $\mathbb{G}$. The
\textbf{lock core} is: $f_1^2>0$ and the atom is isolated ($\Delta>0$ above it). The campaign
summarises this as
\begin{equation}
\text{Clay}\iff\neg(\mathrm{iii}),\qquad
(\mathrm{iii})=\text{``the YM measure is a gapless interacting 4D CFT''.}
\label{eq:lockcore}
\end{equation}

This rewriting is useful because it names the exact analytic object (an atom weight in a positive
measure at the bottom of a channel) and the exact analytic properties it has:
\begin{itemize}
\item $f_1^2$ is \textbf{non-diagonal}: it is an off-diagonal matrix element squared, so no
  diagonal (spectral-density) argument can see it directly;
\item $f_1^2$ is \textbf{upper-semicontinuous} in the coupling: $\nu(\{x\})$ is usc, so
  continuity can only make an atom \emph{fall}, never rise;
\item certifying \textbf{isolation} requires a second scale $E_2$ (the next state), which is why
  every bottom-functional tool is blind;
\item the \textbf{susceptibility} $\chi=\int s^{-1}d\nu$ is a residue-weighted harmonic mean,
  dominated by the lightest residue rather than the lightest mass; this is the technical heart of
  Chapter~\ref{sec:squeeze}.
\end{itemize}

\textbf{Honesty checkpoint.} ``Clay $\iff$ atom'' is \emph{not} a theorem of equivalence. The
campaign's own adversarial audit (\code{research/unconditional/REDUCTION-AUDIT.md}) establishes
that the correct logical shape is a \textbf{diamond}, not a chain
(Figure~\ref{fig:diamond}): a gapped theory need not have an atom at all (the massive free scalar
with $O={:}\varphi^2{:}$ has a gap and $f_1^2=0$), so $\text{gap}\Rightarrow\text{atom}$ is false;
the true direction is \emph{isolated atom at $\inf\operatorname{supp}\nu$ $\Rightarrow$
exponential clustering in that channel}, conditional on that channel carrying the lightest
state. This paper therefore treats ``the atom $f_1^2>0$'' as the \textbf{target object}, not as a
synonym for Clay. The gap between the two is itself a named missing step --- the ordering input
that the lightest state lies in $0^{++}$ (the corpus's \code{IN-11}, a physical input, not
derivable from symmetry).

This edition adds a machine-checked sharpening of the same point. The module
\code{lean/YMStatement/H1Sharpening.lean} proves, with zero \code{sorry} and only the standard
axioms, that the susceptibility ceiling H1 together with the squeeze hypothesis $f_1^2\le\chi m$
and $m>0$ does \emph{not} imply $f_1^2>0$: the tuple $(f,m,M,\chi,\lambda)=(0,1,1,1,1)$ is a
model. The statement is not a numerical accident but a formal theorem
(\code{H1\_does\_not\_imply\_atom}); see the countermodel in \S\ref{thm:leanblind} and
Appendix~\ref{app:lean}.

\figreduction

\subsection{The two halves}

The campaign splits Clay into two halves and a channel:

\begin{itemize}
\item \textbf{G1 (existence).} 4D continuum existence with OS axioms. The fixed-volume
  construction is essentially complete; the one genuinely open UV input is a \textbf{uniform
  renormalization constant} $Z_a$, equivalently a cutoff-uniform bound on the force Hessian (the
  node called L3). L3 is purely ultraviolet and \emph{orthogonal} to the gap --- proving it would
  make fixed-volume existence unconditional but would not help G2.
\item \textbf{G2 (gap).} The proof or transport of $f_1^2>0$ (equivalently of the coarse curvature
  $G''(0)>0$ at $k\to0$). This is where the problem lives.
\item \textbf{The physical channel.} Identifying the lightest state as $0^{++}$ and showing it
  does not dissolve (the no-dissolution theorem). This half is, conditionally, complete.
\end{itemize}

The single deepest structural finding of the campaign is that \textbf{G1's L3 and G2's H1 are one
species} (Figure~\ref{fig:scales}): both are a one-sided \emph{upper bound on a variance}. L3 is
the UV force-variance bound; H1 is the IR conditioned-variance bound
$\operatorname{Var}(\xi_{k\to0})\le1/\lambda$ (curvature form
$\operatorname{Var}_\mu(\langle v,\nabla S\rangle|0)\le E_\mu[H_{vv}|0]+W_{\mathrm{full}}-\delta$).
The shapes are identical because $\operatorname{Var}\ge0$ can only supply a lower edge, while in
the exact identity $G''=E[H]-\operatorname{Var}+W_{\mathrm{full}}$ (with
$W_{\mathrm{full}}=W_{\mathrm{geo}}+X$ as in \S\ref{sec:wgeo}) the variance enters with a minus
sign. The two faces differ only in scale (UV versus IR). This unification is stated in the corpus
as a \textbf{structural identification and a conjecture}, not a theorem
(\code{G2-ASSEMBLE-100.md} \S7: ``a structural identification \ldots{} one closure would likely
close both faces --- but that is a conjecture, not a theorem''). This paper preserves that label.

\figscales

\subsection{The wall in two coordinates (new in this edition)}
\label{sec:preview-twocoord}

The first edition drew the wall in one coordinate: the measure/spectral coordinate, in which the
obstruction is that the atom is a weak-$*$ discontinuity and every finitely computable datum is
weak-$*$ continuous. This edition adds a second coordinate, developed independently and with a
disjoint toolbox: the $\beta$-function / renormalization-group / conformal coordinate, in which
the obstruction is that pure Yang--Mills being gapped is equivalent to $N_f=0$ lying outside the
conformal window, and that statement is dynamical rather than kinematic.

The two coordinates are the subject of Part~III. The short version is this. In the measure
coordinate (Chapter~\ref{ch:weakstar}) the atom is exactly the failure of weak-$*$ continuity of
the boundary density; eight falsification routes fail to build a functional that forces the atom
while being neither a reading (A3) nor an infinite-data limit (H1). In the $\beta$ coordinate
(Chapter~\ref{ch:betarg}) the gap is the absence of a strong-coupling infrared fixed point; the
statement ``a gapless interacting 4D asymptotically-free Wilsonian IR attractor is conformal'' is
true but reduces the problem to that dynamical fact, and eight routes through the modern
$\beta$/CFT toolbox fail, each for a distinct reason. Chapter~\ref{ch:twocoord} assembles the
two-coordinate picture and states what it does and does not prove.

The value of the second coordinate is diagnostic. It shows that the wall is not an artifact of
the spectral language: a completely different set of techniques --- conformal bootstrap, the
$a$-theorem, central charges, large-$N$, holography, Dyson--Schwinger equations, lattice
step-scaling --- hits the same invariant. It also shows, sharply, why each such technique is
insufficient: the bootstrap tests CFT \emph{existence}, not RG \emph{flow}; the $a$-theorem is a
ceiling; central charges are phase-blind; the holographic null-energy condition is compatible
with a gapless infrared throat; and the Dyson--Schwinger scaling solution is gauge-fixed and
unphysical. None of these is a defect of the technique. Each is a precise statement about what
the gap is not.

\figtwocoord

\subsection{What this paper is, and is not}

This paper is \textbf{not} a proof of Clay, and it contains no step labelled QED. It is:

\begin{itemize}
\item a precise statement of the reduction, with every hypothesis displayed;
\item a complete inventory of what was actually proved, labelled by type;
\item a statement of the no-gos, each with the precise reason its class fails;
\item the immune map: on the order of a hundred and fifty routes, each with the exact point at
  which it dies;
\item the two-coordinate map of the wall, with the technique pool that fails in each coordinate
  and the reason it fails;
\item a specification of the missing tool: five conditions, verbatim;
\item an honest verdict.
\end{itemize}

Everything new in this printing is \emph{classification and erratum}: no hypothesis is removed and
no inequality is proved, and \textbf{Clay stays at $0\%$}. The three new primary files are
\code{ROUND1-VERDICTS.md}, \code{RB-RATIO-BOUND.md} and \code{WA-ALIGNMENT.md}, under
\code{research/v3/species/}.

The reader who wants only the bottom line should read Chapter~\ref{sec:verdict}. The reader who
wants to check the work should follow the file names in Chapter~\ref{sec:theorems} and
Appendix~\ref{app:verif}.

\subsection{Relation to the first edition, and to the literature}

Version 1 of this paper (31 pages, DOI 10.5281/zenodo.24486316) reported the reduction, the
ledger, the six tools, the no-gos, the immune map and the verdict. This edition preserves all of
that, expands the proofs, and adds the results of the waves that followed: the falsification wave
(ONLY-EXIT), the tool wave, the Savvidy cleanup, the $\kappa$-line closure, the diagnostics
(P3 approach law, Re-slope, finite-size law), the star wave ($(\star)$), and the $(G1)$ wave. It
also records one numerical erratum (the central-charge normalization, Appendix~\ref{app:errata}).

The paper is deliberately not a survey of the mass-gap literature. The literature is enormous
(constructive field theory, lattice gauge theory, holography, the conformal bootstrap), and the
campaign's contribution is not a new review but a new \emph{audit}: a labelled inventory of what
can and cannot be concluded, with a machine-checked core and a specification of the missing
object. Where the paper touches standard results --- Balaban's renormalization-group programme,
Magnen--Rivasseau--S\'en\'eor's infrared construction, Shen--Zhu--Zhu's strong-coupling gap,
Osterwalder--Seiler's polymer expansion, Banks--Zaks' conformal window, Polchinski's and
Dymarsky--Komargodski--Schwimmer--Theisen's scale-versus-conformal results,
Komargodski--Schwimmer's $a$-theorem --- it cites them and says exactly what it imports.
Where it relies on a numerical measurement, it labels the statement \LE{} and gives the
statistical caveat. Where a statement is a conjecture, it is called a conjecture.

\subsection{Roadmap}

Chapter~\ref{sec:squeeze} proves the exact spectral squeeze and explains why H1 is necessary but
not sufficient, including the machine-checked countermodel. Chapter~\ref{ch:existence} states G1
and the node board. Chapter~\ref{ch:gap} states G2, the region coverage, and the anatomy of the
Savvidy cell. Chapter~\ref{sec:theorems} states the proven theorems with their hypotheses.
Chapter~\ref{ch:tools} presents the tools and each one's free-field acid test, including the tool
wave and its structural classification. Chapter~\ref{sec:nogos} records the no-gos.
Chapter~\ref{ch:immune} states the one-species reduction and the immune map.
Chapter~\ref{ch:weakstar} develops the measure/weak-$*$ coordinate.
Chapter~\ref{ch:betarg} develops the $\beta$/RG/conformal coordinate.
Chapter~\ref{ch:twocoord} assembles the two-coordinate picture.
Chapter~\ref{ch:numerics} reports the numerics and the diagnostics.
Chapter~\ref{sec:missing} specifies what the missing tool must satisfy.
Chapter~\ref{sec:verdict} gives the verdict and the AI contribution statement.

Chapter~\ref{ch:uvface} classifies the ultraviolet face: the last inequality $(\star)$, its exact
one-plaquette counterexample, the good-set theorem, and the periodic constant-mode anchor.

The new chapter on the sixth species (\code{ch16-sixth-species.tex}) reports the outcome of that
wave: \emph{five species classes are closed, the meta-theorem $F3\wedge F4=\varnothing$ is made
quantitative, and two load-bearing constants are corrected} --- and it states explicitly that it
opens no new route.
The search for that sixth species ran on five new classes and repeated the same ending; the detail
is left to that chapter.
Appendix~\ref{app:lean} records the Lean formalization, Appendix~\ref{app:protocols} the
numerical protocols, Appendix~\ref{app:verif} the verification map, and
Appendix~\ref{app:errata} the errata and the compendium of definitions.
